\begin{ex}
Tìm $\frac{2}{3}$ của 15.
\shortans{10}
\loigiai{
$15 \times \frac{2}{3} = 10$.
}
\end{ex}

\begin{ex}
Tìm $\frac{3}{4}$ của 24.
\shortans{18}
\loigiai{
$24 \times \frac{3}{4} = 18$.
}
\end{ex}

\begin{ex}
Tìm $\frac{2}{5}$ của 30.
\shortans{12}
\loigiai{
$30 \times \frac{2}{5} = 12$.
}
\end{ex}

\begin{ex}
Tìm $\frac{5}{6}$ của 42.
\shortans{35}
\loigiai{
$42 \times \frac{5}{6} = 35$.
}
\end{ex}

\begin{ex}
Tìm $\frac{4}{7}$ của 49.
\shortans{28}
\loigiai{
$49 \times \frac{4}{7} = 28$.
}
\end{ex}

\begin{ex}
Tìm $\frac{3}{8}$ của 64.
\shortans{24}
\loigiai{
$64 \times \frac{3}{8} = 24$.
}
\end{ex}

\begin{ex}
Tìm $\frac{5}{9}$ của 81.
\shortans{45}
\loigiai{
$81 \times \frac{5}{9} = 45$.
}
\end{ex}

\begin{ex}
Một rổ có 36 quả cam. Mẹ lấy đi $\frac{2}{3}$ số cam để làm nước ép. Mẹ đã lấy bao nhiêu quả cam?
\shortans{24}
\loigiai{
$36 \times \frac{2}{3} = 24$ quả.
}
\end{ex}

\begin{ex}
Một lớp học có 40 học sinh, trong đó $\frac{3}{5}$ số học sinh là học sinh nữ. Hỏi lớp đó có bao nhiêu học sinh nữ?
\shortans{24}
\loigiai{
$40 \times \frac{3}{5} = 24$ bạn nữ.
}
\end{ex}

\begin{ex}
Lớp đó có bao nhiêu học sinh nam trong bài toán trên?
\shortans{16}
\loigiai{
$40 - 24 = 16$ bạn nam.
}
\end{ex}

\begin{ex}
Một đàn gà có 45 con, trong đó $\frac{4}{9}$ là gà trống. Hỏi đàn gà có bao nhiêu con gà mái?
\shortans{25}
\loigiai{
Gà trống: $45 \times \frac{4}{9} = 20$ con. Gà mái: $45 - 20 = 25$ con.
}
\end{ex}

\begin{ex}
Một mảnh đất hình chữ nhật dài 60 m, chiều rộng bằng $\frac{3}{4}$ chiều dài. Chiều rộng mảnh đất là bao nhiêu mét?
\shortans{45}
\loigiai{
$60 \times \frac{3}{4} = 45$ m.
}
\end{ex}

\begin{ex}
Diện tích mảnh đất trong bài toán trên là bao nhiêu mét vuông?
\shortans{2700}
\loigiai{
$60 \times 45 = 2700\text{ m}^2$.
}
\end{ex}

\begin{ex}
Một cuốn sách có 120 trang, Lan đã đọc được $\frac{5}{6}$ số trang. Lan đã đọc được bao nhiêu trang sách?
\shortans{100}
\loigiai{
$120 \times \frac{5}{6} = 100$ trang.
}
\end{ex}

\begin{ex}
Số trang sách Lan chưa đọc trong bài toán trên là bao nhiêu trang?
\shortans{20}
\loigiai{
$120 - 100 = 20$ trang.
}
\end{ex}

\begin{ex}
Một người có 100 000 đồng, người đó đã tiêu hết $\frac{3}{5}$ số tiền đó. Người đó đã tiêu bao nhiêu nghìn đồng?
\shortans{60}
\loigiai{
$100 \times \frac{3}{5} = 60$ nghìn đồng.
}
\end{ex}

\begin{ex}
Một thùng có 75 lít dầu, người ta lấy ra $\frac{2}{5}$ số dầu. Hỏi đã lấy ra bao nhiêu lít dầu?
\shortans{30}
\loigiai{
$75 \times \frac{2}{5} = 30$ lít.
}
\end{ex}

\begin{ex}
Trong thùng còn lại bao nhiêu lít dầu trong bài toán trên?
\shortans{45}
\loigiai{
$75 - 30 = 45$ lít.
}
\end{ex}

\begin{ex}
Tìm $\frac{7}{10}$ của 1000.
\shortans{700}
\loigiai{
$1000 \times \frac{7}{10} = 700$.
}
\end{ex}

\begin{ex}
Tìm $\frac{3}{5}$ của 250 kg theo đơn vị ki-lô-gam.
\shortans{150}
\loigiai{
$250 \times \frac{3}{5} = 150$ kg.
}
\end{ex}

\begin{ex}
Tìm $\frac{5}{12}$ của một năm (năm có 12 tháng) theo số tháng.
\shortans{5}
\loigiai{
$12 \times \frac{5}{12} = 5$ tháng.
}
\end{ex}

\begin{ex}
Tìm $\frac{3}{4}$ của một giờ (giờ có 60 phút) theo số phút.
\shortans{45}
\loigiai{
$60 \times \frac{3}{4} = 45$ phút.
}
\end{ex}

\begin{ex}
Tìm $\frac{2}{3}$ của một ngày (ngày có 24 giờ) theo số giờ.
\shortans{16}
\loigiai{
$24 \times \frac{2}{3} = 16$ giờ.
}
\end{ex}

\begin{ex}
Tìm $\frac{1}{4}$ của một thế kỉ (thế kỉ có 100 năm) theo số năm.
\shortans{25}
\loigiai{
$100 \times \frac{1}{4} = 25$ năm.
}
\end{ex}

\begin{ex}
Biết $\frac{2}{3}$ của một số là 20. Số đó là bao nhiêu?
\shortans{30}
\loigiai{
$20 : \frac{2}{3} = 20 \times \frac{3}{2} = 30$.
}
\end{ex}
